\documentclass[11pt]{article}

\usepackage[T1]{fontenc}
\usepackage{lmodern}
\usepackage{microtype}
\usepackage[a4paper,margin=27mm,headheight=14pt]{geometry}
\usepackage{amsmath,amssymb,amsthm,mathtools}
\usepackage{booktabs,longtable,array}
\usepackage{enumitem}
\usepackage{xcolor}
\usepackage{pid-rs-report-tables}
\usepackage{hyperref}
\usepackage{fancyhdr}

\setlength{\emergencystretch}{2em}
\newcolumntype{L}[1]{>{\raggedright\arraybackslash}p{#1}}

\ifdefined\pdfinfoomitdate
  \pdfinfoomitdate=1
\fi
\ifdefined\pdftrailerid
  \pdftrailerid{}
\fi
\ifdefined\pdfsuppressptexinfo
  \pdfsuppressptexinfo=15
\fi

\hypersetup{
  colorlinks=true,
  linkcolor=blue!55!black,
  citecolor=blue!55!black,
  urlcolor=blue!55!black,
  pdftitle={Support-Change-Tolerant Continuity for Averaged Categorical Shared-Exclusions PID},
  pdfsubject={Closed-simplex total-variation bounds, exact counterexamples, and evidence boundary},
  pdfkeywords={partial information decomposition, shared exclusions, continuity, total variation},
  pdfauthor={pid-rs project analysis}
}

\PidConfigureSections
\PidSetRunningHeads{pid-rs averaged categorical SxPID}{support-change analysis}
\PidConfigureAbstract
\PidConfigureLinksAndListings

\newtheorem{theorem}{Theorem}[section]
\newtheorem{corollary}[theorem]{Corollary}
\newtheorem{lemma}[theorem]{Lemma}
\newtheorem{proposition}[theorem]{Proposition}
\theoremstyle{definition}
\newtheorem{definition}[theorem]{Definition}
\newtheorem{remark}[theorem]{Remark}
\newtheorem{counterexample}[theorem]{Counterexample}

\newcommand{\cZ}{\mathcal{Z}}
\newcommand{\cS}{\mathcal{S}}
\newcommand{\supp}{\operatorname{supp}}
\newcommand{\TV}{d_{\mathrm{TV}}}
\newcommand{\normone}[1]{\left\lVert #1\right\rVert_1}
\newcommand{\sx}{\mathrm{sx}}
\newcommand{\Ind}{\mathbf{1}}

\title{\PidReportTitle
  {Support-Change-Tolerant Continuity\\
   for Averaged Categorical Shared-Exclusions PID}
  {Closed-simplex total-variation bounds, exact counterexamples,\\
   and a finite-sample transfer without a positive support-mass floor}}
\author{pid-rs project analysis}
\date{24 July 2026}

\begin{document}
\PidMakeTitle

\begin{abstract}
Categorical shared-exclusions partial information decomposition (SxPID) is defined pointwise and
then averaged with the same joint law.  Existing interior differentiability and fixed-support
arguments do not by themselves control this average when cells enter or leave the support.
This document gives a project-defined exact-real continuity theorem on a fixed finite ambient
alphabet.  A pointwise overlap decomposition separates residual entropy from a common-overlap
neighborhood load.  For an event that is a union of \(J\) equivalence neighborhoods, the load is at
most \(J\) times total variation and yields the modulus
\[
  (1-\eta)\log\!\left(1+\frac{J\eta}{1-\eta}\right).
\]
The construction is sharp for this decomposition and is realizable by a paper-semantic Sx event.
Published nonnegativity of the informative and misinformative pointwise atoms sharpens the
full-lattice M{\"o}bius transfer: component residuals use a maximum of two residual entropies,
whereas signed net residuals require their sum.  Finite-alphabet entropy envelopes produce an
explicit modulus that vanishes without a positive population support-mass floor.  Monotone
relaxations compose with a dependency-colored empirical-law radius.  Exact categorical
counterexamples reject a global linear modulus, pointwise boundary continuity, an active-face
Fannes--Audenaert substitution, a maximum-residual net shortcut, and alphabet-free control.
Lean checks the exact heterogeneous keyed Sx events, the finite equivalence-union load bound, and,
separately, supplied exact natural counts with positive total through the four fixed two-source
signed-net cumulative logarithms and positive-support averages.  It does not check the
support-change logarithmic transfer between laws, complete informative or misinformative averaged
components, concrete atoms, bytes-to-counts, Rust or floating-point refinement, certifier/parser
execution, or statistical and population validity.
\end{abstract}

\begingroup
\small
\setlength{\parskip}{0pt}
\tableofcontents
\endgroup

\section{Status, scope, and provenance}

The categorical Sx functional, redundancy lattice, and pointwise component-sign theorem are due to
Makkeh, Gutknecht, and Wibral~\cite{makkeh2021,gutknecht2021}.  This document analyzes that
paper-defined functional.  The support-change theorem and its validation artifacts are
project-defined.  They do not define another PID measure and make no scientific-priority claim.

Fix nonempty finite source alphabets, a nonempty finite target alphabet, the complete
Cartesian-product alphabet, and the full finite redundancy lattice.  These objects and the event map remain fixed while the
law changes.  All logarithms are natural and all mathematical statements use exact real
arithmetic.
\begin{equation}
  \cZ
  =
  \cS_1\times\cdots\times\cS_m\times\mathcal T,
  \qquad
  K=|\cZ|.
\end{equation}

The result applies to averaged categorical cumulatives and atoms.  It does not apply to:
\begin{itemize}[leftmargin=2em]
  \item a changing observed alphabet or selected lattice;
  \item pointwise values at a key whose mass vanishes;
  \item a same-sample fitted or changing quantizer;
  \item continuous, mixed, or singular shared exclusions;
  \item portable binary64 signs or correctly rounded logarithms; or
  \item statistical or consumer validity without a separate contract.
\end{itemize}

\section{Overlap and residuals}

Let \(p,q\) be probability laws on one finite set \(X\).  Put
\begin{equation}
  \eta=\TV(p,q)=\frac12\normone{p-q},
\end{equation}
and define
\begin{equation}
  r_x=\min\{p_x,q_x\},
  \qquad
  a_x=p_x-r_x,
  \qquad
  b_x=q_x-r_x.
\end{equation}
Then
\begin{equation}
  p=r+a,\qquad q=r+b,
\end{equation}
\begin{equation}\label{eq:mass-identities}
  \sum_xa_x=\sum_xb_x=\eta,
  \qquad
  \sum_xr_x=1-\eta=:R,
\end{equation}
and \(a_xb_x=0\) for every \(x\).
For any law functional \(F\), write
\begin{equation}
  \Delta F=F(p)-F(q).
\end{equation}

For a nonnegative finite vector \(d\), define its residual entropy with the continuous zero
convention:
\begin{equation}
  E(d)=-\sum_{x:d_x>0}d_x\log d_x.
\end{equation}
Write
\begin{equation}
  E_\vee=\max\{E(a),E(b)\},
  \qquad
  E_\Sigma=E(a)+E(b).
\end{equation}
These are entropies of subprobability vectors, not entropies of separately normalized laws.

\section{Anchored-neighborhood continuity}

Fix a set \(N_x\subseteq X\) for every \(x\), subject to the anchor condition
\begin{equation}\label{eq:anchor}
  x\in N_x.
\end{equation}
For a nonnegative vector \(v\), write
\begin{equation}
  v(N_x)=\sum_{y\in N_x}v_y.
\end{equation}
Define
\begin{equation}\label{eq:G}
  G_N(p)=-\sum_{x:p_x>0}p_x\log p(N_x).
\end{equation}
Equation~\eqref{eq:anchor} makes every displayed logarithm well-defined.

For \(R>0\), define the overlap load
\begin{equation}\label{eq:T}
  T_N(r,d)
  =
  \sum_{x:r_x>0}
  r_x\frac{d(N_x)}{r(N_x)}.
\end{equation}

\begin{theorem}[Residual-plus-load bound]\label{thm:generic}
For any two probability laws on \(X\),
\begin{equation}\label{eq:generic-bound}
  |G_N(p)-G_N(q)|
  \le
  E_\vee+
  R\log\!\left(
    1+\frac{\max\{T_N(r,a),T_N(r,b)\}}{R}
  \right),
\end{equation}
where the logarithmic term is defined as zero when \(R=0\).
\end{theorem}

\begin{proof}
Substitution gives the exact identity
\begin{equation}\label{eq:exact-split}
  G_N(p)-G_N(q)=C_N+U_{p,a}-U_{q,b},
\end{equation}
where
\begin{equation}
  C_N
  =
  \sum_{x:r_x>0}
  r_x
  \log
  \frac{r(N_x)+b(N_x)}{r(N_x)+a(N_x)},
\end{equation}
and
\begin{equation}
  U_{p,a}=-\sum_{x:a_x>0}a_x\log p(N_x),
  \qquad
  U_{q,b}=-\sum_{x:b_x>0}b_x\log q(N_x).
\end{equation}
The anchor condition gives
\begin{equation}
  a_x\le p(N_x)\le1,
  \qquad
  b_x\le q(N_x)\le1.
\end{equation}
Consequently,
\begin{equation}\label{eq:residual-interval}
  0\le U_{p,a}\le E(a),
  \qquad
  0\le U_{q,b}\le E(b),
\end{equation}
and
\begin{equation}
  |U_{p,a}-U_{q,b}|\le E_\vee.
\end{equation}

For \(d\in\{a,b\}\), let
\begin{equation}
  C_N(d)=
  \sum_{x:r_x>0}
  r_x\log\!\left(1+\frac{d(N_x)}{r(N_x)}\right).
\end{equation}
Then \(C_N=C_N(b)-C_N(a)\).  Both terms are nonnegative, so
\begin{equation}
  |C_N|\le\max\{C_N(a),C_N(b)\}.
\end{equation}
Jensen's inequality with weights \(r_x/R\) gives
\begin{equation}
  C_N(d)
  \le
  R\log\!\left(1+\frac{T_N(r,d)}R\right).
\end{equation}
Combining the bounds proves Equation~\eqref{eq:generic-bound}.  If \(R=0\), the common sum is empty.
If \(\eta=0\), both residuals vanish.  Thus the stated convention covers both endpoints.
\end{proof}

\section{Union-of-equivalence load}

Fix an integer \(J\ge1\) and suppose
\begin{equation}\label{eq:union-equivalence}
  N_x=\bigcup_{j=1}^{J}[x]_j,
\end{equation}
where, for each \(j\), \([x]_j\) is the equivalence class of \(x\) under a fixed relation
\(\sim_j\).  The \(J\) relations may differ.

\begin{lemma}[Fractional-cover load]\label{lem:fractional-cover}
If \(d\) is nonnegative and \(\sum_xd_x=\eta\), then
\begin{equation}
  T_N(r,d)\le J\eta.
\end{equation}
\end{lemma}

\begin{proof}
For \(r_x>0\),
\begin{equation}
  \frac{d(N_x)}{r(N_x)}
  \le
  \sum_{j=1}^{J}\frac{d([x]_j)}{r([x]_j)}.
\end{equation}
Fix \(j\) and group \(x\) by an equivalence class \(D\).  Every class with \(r(D)>0\) contributes
\begin{equation}
  \sum_{x\in D}r_x\frac{d(D)}{r(D)}=d(D).
\end{equation}
The positive-overlap classes contribute at most \(\eta\).  A class with no overlap has no anchored
\(x\) in the defining sum and need not contribute its residual mass.  Summing the \(J\) relations
proves the result.
\end{proof}

Define
\begin{equation}\label{eq:gJ}
  g_J(\eta)=
  \begin{cases}
    (1-\eta)\log\!\left(1+\dfrac{J\eta}{1-\eta}\right),
      &0<\eta<1,\\[6pt]
    0,&\eta\in\{0,1\}.
  \end{cases}
\end{equation}

\begin{corollary}[Equivalence-union bound]\label{cor:gJ}
Under Equation~\eqref{eq:union-equivalence},
\begin{equation}
  |G_N(p)-G_N(q)|\le E_\vee+g_J(\eta).
\end{equation}
Moreover,
\begin{equation}\label{eq:g-linear}
  g_J(\eta)\le J\eta.
\end{equation}
\end{corollary}

\begin{proof}
Apply Lemma~\ref{lem:fractional-cover} in Theorem~\ref{thm:generic}.  The linear envelope follows
from \(\log(1+x)\le x\).
\end{proof}

\begin{remark}[Upper radii]
The exact function \(g_J\) is not monotone on all of \([0,1]\).  A proof that knows only
\(\eta\le\varepsilon\) cannot substitute \(\varepsilon\) blindly into
Equation~\eqref{eq:gJ}.  Equation~\eqref{eq:g-linear} is a safe monotone relaxation.
\end{remark}

\begin{proposition}[Paper-semantic equality witness]\label{prop:sharp}
For every \(J\ge1\) and \(0\le\eta\le1\), there is a categorical Sx event for which
\begin{equation}
  G_N(p)-G_N(q)=E(a)+g_J(\eta).
\end{equation}
\end{proposition}

\begin{proof}
Use \(J\) sources and states
\begin{equation}
  z,\quad c,\quad\ell_1,\ldots,\ell_J.
\end{equation}
For source \(j\), make \(c\) and \(\ell_j\) share one categorical value.  Give every other
state-coordinate pair a distinct value.  The singleton antichain has
\begin{equation}
  N_{\ell_j}=\{\ell_j,c\},
  \qquad
  N_z=\{z\},
  \qquad
  N_c=\{c,\ell_1,\ldots,\ell_J\}.
\end{equation}
Let the target be constant and set
\begin{equation}
  r_{\ell_j}=\frac{1-\eta}{J},
  \qquad
  a_z=\eta,
  \qquad
  b_c=\eta.
\end{equation}
The residual at \(z\) contributes \(E(a)=\eta\log(1/\eta)\).  The residual at \(c\) contributes
zero.  Each leaf contributes its share of \(g_J\).  Direct summation proves the equality.
\end{proof}

\section{Categorical shared-exclusions specialization}

Specialize the generic finite set to \(X=\cZ\).

Let \(\beta\) be a node of the full finite redundancy lattice.  For every ambient key
\(z=(s_1,\ldots,s_m,t)\in\cZ\), define
\begin{equation}
  A_\beta(z)
  =
  \bigcup_{u\in\beta}\{y:y_u=s_u\},
\end{equation}
\begin{equation}
  C(z)=\{y:y_T=t\},
  \qquad
  B_\beta(z)=A_\beta(z)\cap C(z).
\end{equation}
These are the finite-alphabet event semantics in Equations (15a)--(15b) and Appendix B
of Makkeh, Gutknecht, and Wibral~\cite{makkeh2021}.  The notation makes the law-independent keyed
event map explicit; it does not alter the paper-defined functional.
If \(J_\beta=|\beta|\), then \(A_\beta\) and \(B_\beta\) are unions of \(J_\beta\) equivalence
neighborhoods.  The target event is one equivalence neighborhood.
The event map is fixed on the complete alphabet, independently of the law.  Local logarithms and
averages below evaluate these events only at positive-mass keys.

Support-restricted averaging gives
\begin{equation}\label{eq:sx-as-G}
  I_\beta^+(p)=G_{A_\beta}(p),
  \qquad
  I_\beta^-(p)=G_{B_\beta}(p)-G_C(p),
\end{equation}
\begin{equation}
  I_\beta^{\sx}(p)
  =
  G_{A_\beta}(p)-G_{B_\beta}(p)+G_C(p).
\end{equation}

\begin{theorem}[Averaged cumulative bounds]\label{thm:cumulative}
For every full-lattice node \(\beta\),
\begin{align}
  |\Delta I_\beta^+|
  &\le E_\vee+g_{J_\beta}(\eta),\label{eq:c-plus}\\
  |\Delta I_\beta^-|
  &\le E_\vee+g_{J_\beta}(\eta)+g_1(\eta),\label{eq:c-minus}\\
  |\Delta I_\beta^\sx|
  &\le E_\Sigma+2g_{J_\beta}(\eta)+g_1(\eta).\label{eq:c-net}
\end{align}
\end{theorem}

\begin{proof}
For every supported key,
\begin{equation}\label{eq:local-ranges}
  0\le c_\beta^+(z;p)\le-\log p_z,
  \qquad
  0\le c_\beta^-(z;p)\le-\log p_z,
\end{equation}
and
\begin{equation}
  |c_\beta^\sx(z;p)|\le-\log p_z.
\end{equation}
\(z\in B_\beta(z)\subseteq C(z)\) and \(z\in A_\beta(z)\), so each keyed event probability is at
least \(p_z\), while every event probability is at most one.
The informative residual averages therefore differ by at most \(E_\vee\).  Treating the
misinformative expression in Equation~\eqref{eq:sx-as-G} as one local value gives the same
\(E_\vee\), not two residual entropies.  Its common term has one \(B_\beta\) and one \(C\)
contribution.  The signed net residual averages require \(E(a)+E(b)\); its common term has
\(A_\beta\), \(B_\beta\), and \(C\) contributions.  Apply
Corollary~\ref{cor:gJ} to each contribution.
\end{proof}

\section{Full-lattice atom transfer}

Let \(M\) be the fixed M{\"o}bius inverse in the orientation
\begin{equation}
  \Pi_\alpha^u=\sum_\beta M_{\alpha\beta}I_\beta^u.
\end{equation}
Define
\begin{equation}
  W_\alpha(\eta)=
  \sum_\beta|M_{\alpha\beta}|g_{J_\beta}(\eta),
\end{equation}
and
\begin{equation}
  s_\alpha=\sum_\beta M_{\alpha\beta}.
\end{equation}

Let \(\bot\) be the unique least node and \(Z\) the down-set zeta matrix.  Since
\begin{equation}
  Z e_\bot=\mathbf1,
\end{equation}
we have
\begin{equation}\label{eq:row-sum}
  M\mathbf1=e_\bot,
  \qquad
  s_\alpha=\Ind\{\alpha=\bot\}.
\end{equation}

Equation (19), M{\"o}bius inversion, and Theorem IV.3 of the defining paper~\cite{makkeh2021}
give
\begin{equation}\label{eq:component-sign}
  \pi_\alpha^+(z;p)\ge0,
  \qquad
  \pi_\alpha^-(z;p)\ge0.
\end{equation}
They reconstruct the greatest-node component cumulative.  At that node, for a supported key
\(z=(s,t)\),
\begin{equation}
  c_\top^+(z;p)=-\log p(S=s)\le-\log p_z,
  \qquad
  c_\top^-(z;p)=-\log p(S=s\mid T=t)\le-\log p_z.
\end{equation}
Consequently, for every atom,
\begin{equation}\label{eq:atom-surprisal}
  0\le\pi_\alpha^+(z;p)\le-\log p_z,
  \qquad
  0\le\pi_\alpha^-(z;p)\le-\log p_z.
\end{equation}

\begin{theorem}[Averaged atom bounds]\label{thm:atoms}
For every atom \(\alpha\) of the full finite redundancy lattice,
\begin{align}
  |\Delta\Pi_\alpha^+|
  &\le E_\vee+W_\alpha(\eta),\label{eq:a-plus}\\
  |\Delta\Pi_\alpha^-|
  &\le E_\vee+W_\alpha(\eta)+|s_\alpha|g_1(\eta),\label{eq:a-minus}\\
  |\Delta\Pi_\alpha^\sx|
  &\le E_\Sigma+2W_\alpha(\eta)+|s_\alpha|g_1(\eta).
  \label{eq:a-net}
\end{align}
\end{theorem}

\begin{proof}
Split each averaged atom into overlap, left-residual, and right-residual weighted pointwise atoms.
For the overlap part, commute the fixed finite M{\"o}bius matrix with the finite sum.  The
\(A_\beta\) and \(B_\beta\) common terms receive the row coefficients \(M_{\alpha\beta}\).
Equation~\eqref{eq:row-sum} gives the coefficient of the common target term.

For a component atom, Equations~\eqref{eq:component-sign} and~\eqref{eq:atom-surprisal} place each
residual average in \([0,E(a)]\) or \([0,E(b)]\).  Their difference is at most \(E_\vee\).
For a net atom,
\begin{equation}
  \pi_\alpha^\sx=\pi_\alpha^+-\pi_\alpha^-,
\end{equation}
so
\begin{equation}
  |\pi_\alpha^\sx(z;p)|\le-\log p_z.
\end{equation}
The two signed residual averages require \(E_\Sigma\).  The common-term bounds give the remaining
terms in Equations~\eqref{eq:a-plus}--\eqref{eq:a-net}.
\end{proof}

\begin{remark}[Why the full lattice matters]
Equation~\eqref{eq:row-sum} requires a least node.  The sign theorem in
Equation~\eqref{eq:component-sign} also concerns the paper-defined full lattice.  Neither fact
transfers automatically to a truncated node family or an arbitrary coefficient matrix.
\end{remark}

\section{Finite-alphabet envelopes and uniform continuity}

Assume \(\eta>0\), and put
\begin{equation}
  k_a=|\supp a|,
  \qquad
  k_b=|\supp b|.
\end{equation}
The residual supports are disjoint, so \(k_a+k_b\le K\).  Concavity of \(-x\log x\) gives
\begin{equation}
  E(a)\le\eta\log\frac{k_a}{\eta},
  \qquad
  E(b)\le\eta\log\frac{k_b}{\eta}.
\end{equation}
Therefore
\begin{equation}\label{eq:e-max}
  E_\vee
  \le
  e_K^\vee(\eta)
  :=
  \eta\log\frac{K-1}{\eta},
\end{equation}
and
\begin{equation}\label{eq:e-sum}
  E_\Sigma
  \le
  e_K^\Sigma(\eta)
  :=
  \eta\log\frac{\lfloor K^2/4\rfloor}{\eta^2}.
\end{equation}
Define both envelopes as zero at \(\eta=0\).  If \(K=1\), then \(\eta=0\).

\begin{corollary}[Closed-simplex continuity]\label{cor:continuity}
On a fixed finite Cartesian-product alphabet and full finite redundancy lattice, every averaged
categorical shared-exclusions cumulative and atom is uniformly continuous in total variation on
the closed probability simplex.  The result remains valid when cells enter or leave the support
and requires no positive minimum cell mass.
\end{corollary}

\begin{proof}
Substitute Equations~\eqref{eq:e-max} and~\eqref{eq:e-sum} into
Theorems~\ref{thm:cumulative} and~\ref{thm:atoms}.  Every resulting term tends to zero as
\(\eta\downarrow0\).  The finite family of nodes can be bounded simultaneously by taking its
maximum.
\end{proof}

\subsection{Range caps}

Let
\begin{equation}
  K_S=\prod_{i=1}^{m}|\mathcal S_i|.
\end{equation}
The local ranges give
\begin{equation}
  0\le I_\beta^+(p)\le H_p(S)\le\log K_S,
\end{equation}
\begin{equation}
  0\le I_\beta^-(p)\le H_p(S\mid T)\le\log K_S.
\end{equation}
The same upper bounds hold for every averaged component atom because nonnegative component atoms
sum to the greatest-node component cumulative.  Hence every component difference can be capped
by \(\log K_S\), and every net difference can be capped by \(2\log K_S\).

\section{Support-change-tolerant statistical transfer}

The exact functions \(g_J,e_K^\vee,e_K^\Sigma\) are not all monotone on \([0,1]\).  To use only an
upper radius \(0\le\eta\le\varepsilon\le1\), define
\begin{equation}
  \bar e_K^\vee(\varepsilon)
  =
  \varepsilon
  \left(1+\log\frac{K-1}{\varepsilon}\right),
\end{equation}
\begin{equation}
  \bar e_K^\Sigma(\varepsilon)
  =
  \varepsilon
  \left(
    2+\log\frac{\lfloor K^2/4\rfloor}{\varepsilon^2}
  \right),
\end{equation}
with value zero at \(\varepsilon=0\).  These are monotone on \([0,1]\), and
\begin{equation}
  E_\vee\le\bar e_K^\vee(\varepsilon),
  \qquad
  E_\Sigma\le\bar e_K^\Sigma(\varepsilon),
  \qquad
  g_J(\eta)\le J\varepsilon.
\end{equation}

To verify the first two inequalities, put
\(A=K-1\ge1\) and \(B=\lfloor K^2/4\rfloor\ge1\).  For
\(0<\eta\le\varepsilon\), both \(\log(A/\varepsilon)\) and
\(\log(B/\varepsilon^2)\) are nonnegative, while
\begin{equation}
  \eta\log\frac{\varepsilon}{\eta}
  =
  \varepsilon
  \left(\frac{\eta}{\varepsilon}\right)
  \log\left(\frac{\varepsilon}{\eta}\right)
  \le\frac{\varepsilon}{e}
  \le\varepsilon.
\end{equation}
Therefore
\begin{equation}
  \eta\log\frac{A}{\eta}
  \le
  \varepsilon\log\frac{A}{\varepsilon}+\varepsilon,
\end{equation}
and
\begin{equation}
  \eta\log\frac{B}{\eta^2}
  \le
  \varepsilon\log\frac{B}{\varepsilon^2}+2\varepsilon.
\end{equation}
The endpoint \(\eta=0\) follows by the zero convention.  The last inequality follows from
\(g_J(\eta)\le J\eta\le J\varepsilon\).

Put
\begin{equation}
  L_\alpha=\sum_\beta|M_{\alpha\beta}|J_\beta.
\end{equation}
On any event where \(\TV(\widehat p,p)\le\varepsilon\), Table~\ref{tab:transfer} holds.

\begin{table}[ht]
\centering
\caption{Monotone support-change transfer from a total-variation radius.}
\label{tab:transfer}
\begin{tabular}{L{0.32\textwidth}L{0.58\textwidth}}
\toprule
\PidTableHeaderRow
Quantity & Radius \\
\midrule
\(I_\beta^+\) &
\(\bar e_K^\vee(\varepsilon)+J_\beta\varepsilon\) \\
\(I_\beta^-\) &
\(\bar e_K^\vee(\varepsilon)+(J_\beta+1)\varepsilon\) \\
\(I_\beta^\sx\) &
\(\bar e_K^\Sigma(\varepsilon)+(2J_\beta+1)\varepsilon\) \\
\(\Pi_\alpha^+\) &
\(\bar e_K^\vee(\varepsilon)+L_\alpha\varepsilon\) \\
\(\Pi_\alpha^-\) &
\(\bar e_K^\vee(\varepsilon)+(L_\alpha+|s_\alpha|)\varepsilon\) \\
\(\Pi_\alpha^\sx\) &
\(\bar e_K^\Sigma(\varepsilon)+(2L_\alpha+|s_\alpha|)\varepsilon\) \\
\bottomrule
\end{tabular}
\end{table}

The dependency-colored empirical-law theorem in the companion pid-rs analysis supplies an
\(L^1\) radius \(D_n\) under a declared mutual-within-color independence contract.  Set
\begin{equation}
  \varepsilon_n=D_n/2.
\end{equation}
Table~\ref{tab:transfer} then gives an averaged categorical SxPID radius without a population
support-mass floor.  The composition does not validate the coloring from data and does not remove
the exponential alphabet factor in the law-distance theorem.

For a nonidentical-law theorem centered at \(\bar p_n\), let
\begin{equation}
  b_n=\normone{\bar p_n-p_\star}.
\end{equation}
If \(D_n\) controls \(\normone{\widehat p_n-\bar p_n}\), use
\begin{equation}
  \varepsilon_n
  =
  \min\left\{1,\frac{D_n+b_n}{2}\right\}.
\end{equation}
An unknown \(b_n\) remains an unknown scientific bias.

\section{Exact negative results}

\begin{counterexample}[No global linear or pointwise boundary modulus]
Let \(S_2\) be constant and \(T=S_1\).  Give the rare copied state probability \(\eta\) and compare
with the point mass at \(\eta=0\).  The averaged unique-\(S_1\) informative and net atom is
\begin{equation}
  h_2(\eta)
  =
  -\eta\log\eta-(1-\eta)\log(1-\eta).
\end{equation}
The two-source bottom event is the whole sample space because the \(S_2\) branch is constant, so
its local cumulative is zero.  M{\"o}bius inversion therefore makes the unique-\(S_1\) atom equal
to the \(S_1\)-node cumulative.  The node's net local value is pointwise mutual information, which
reduces to source surprisal when \(T=S_1\); its averaged cumulative is \(I(S_1;T)\).
The \(L^1\) distance is \(2\eta\), and
\begin{equation}
  \frac{h_2(\eta)}{2\eta}\longrightarrow\infty.
\end{equation}
The rare-key pointwise atom is \(-\log\eta\).  Therefore no global linear modulus exists, the
\(\eta\log(1/\eta)\) order is necessary, and pointwise support-boundary continuity is false.
\end{counterexample}

\begin{counterexample}[An active-face Fannes--Audenaert substitution fails]
Use three sources, a constant target, the singleton antichain, and source states
\begin{equation}
  000,\quad011,\quad202,\quad330,\quad444.
\end{equation}
Set
\begin{equation}
  p=(0,3/10,3/10,3/10,1/10),
\end{equation}
\begin{equation}
  q=(1/10,3/10,3/10,3/10,0).
\end{equation}
Then \(\eta=1/10\) and
\begin{equation}
  |G_N(p)-G_N(q)|
  =
  \frac1{10}\log10+\frac9{10}\log\frac43.
\end{equation}
The union-support face has
\begin{equation}
  K_\cup
  =
  |\operatorname{supp}p\cup\operatorname{supp}q|
  =5.
\end{equation}
The ordinary Fannes--Audenaert entropy radius on that five-state face is
\begin{equation}
  h_2(1/10)+\frac1{10}\log4.
\end{equation}
It is smaller.  The exact excess is
\begin{equation}
  \frac9{10}\log\frac65-\frac1{10}\log4>0,
\end{equation}
equivalently
\begin{equation}
  6^9>4\cdot5^9.
\end{equation}
This is an actual categorical Sx event.  It rejects substituting the ordinary entropy continuity
radius computed on the union-support face for the neighborhood-functional bound, and it attains
the residual-plus-\(g_3\) bound.  It does not say that the same expression formed with a larger
ambient Cartesian-product cardinality must fail.  A bounded search found no four-state example;
that non-finding does not prove global minimality.
\end{counterexample}

\begin{counterexample}[The signed residual maximum shortcut fails]
Keep \(S_2\) constant.  For \(0<\eta<1\) and \(R=1-\eta\), put
\begin{align}
  p(0,0)&=\eta,&
  p(1,2)&=p(2,1)=R/2,\\
  q(1,1)&=\eta,&
  q(1,2)&=q(2,1)=R/2.
\end{align}
The displayed pairs are \((S_1,T)\).  As above, the constant-\(S_2\) branch makes the bottom net
cumulative zero, so the two-source M{\"o}bius inversion identifies the unique-\(S_1\) net atom
with the \(S_1\)-node cumulative.  The node's local value is pointwise mutual information, while
its averaged cumulative is \(I(S_1;T)\).  The difference of the two residual-weighted
contributions is
\begin{equation}
  2\eta\log\frac{1+\eta}{2\eta}
  >
  \eta\log\frac1\eta.
\end{equation}
The strict inequality reduces to \((1-\eta)^2>0\).  Thus the signed residual budgets cannot be
replaced by their maximum in this proof.  The example does not prove global sharpness of the whole
net-atom modulus.
\end{counterexample}

\begin{counterexample}[A truncated-family row identity fails]
For a V-shaped three-element poset with two incomparable minimal nodes and one greatest node,
constant cumulatives \(c=(1,1,1)\) invert to
\begin{equation}
  \pi=(1,1,-1).
\end{equation}
The greatest-row M{\"o}bius sum is \(-1\).  The full-lattice row-sum and component-sign arguments
cannot be transferred to an arbitrary truncated family.
\end{counterexample}

\begin{counterexample}[No alphabet-free modulus]
Move total mass \(\eta\) from one cell and spread it uniformly over \(M\) new cells in a copied
source construction.  The averaged unique component contains \(\eta\log M\).  A modulus
independent of the fixed ambient alphabet cannot hold.
\end{counterexample}

\begin{proposition}[Worst-case leading logarithmic coefficients]\label{prop:leading-sharp}
There exist fixed finite categorical shared-exclusions systems and fixed atom coordinates whose
component and net differences have leading coefficients one and two, respectively, in units of
\(\eta\log(1/\eta)\) as \(\eta\downarrow0\).  Consequently, consider a family that assigns every
covered fixed system \(\mathcal F\) a bound
\(c\eta\log(1/\eta)+O_{\mathcal F}(\eta)\), where the lower-order constant may depend on
\(\mathcal F\) but \(c\) is common to the family.  Such a family requires \(c\geq1\) for
components and \(c\geq2\) for net atoms.
\end{proposition}

\begin{proof}
For fixed \(K\geq2\), the upper entropy terms satisfy
\begin{equation}
  e_K^\vee(\eta)
  =
  \eta\log\frac1\eta+O(\eta),
  \qquad
  e_K^\Sigma(\eta)
  =
  2\eta\log\frac1\eta+O(\eta),
\end{equation}
and \(g_J(\eta)=O(\eta)\) for fixed \(J\).  Hence every \(W_\alpha\) and row-sum-weighted
\(g_1\) term is \(O(\eta)\) when the finite lattice, M{\"o}bius matrix, and branch counts are
fixed.

The copied-source path gives
\begin{equation}
  \lim_{\eta\downarrow0}
  \frac{h_2(\eta)}{\eta\log(1/\eta)}
  =1.
\end{equation}
This gives one fixed system and atom with informative-component coefficient one.  The
constant-target equality witness has identical informative and misinformative bottom atoms and
gives one fixed-system misinformative witness with the same coefficient.

For the fixed two-source system in the signed-residual counterexample, the complete
unique-\(S_1\) net atom is \(I(S_1;T)\).  With \(R=1-\eta\), direct evaluation gives
\begin{equation}
  I_p(S_1;T)
  =
  -\eta\log\eta-R\log\frac{R}{2},
\end{equation}
and
\begin{equation}
  I_q(S_1;T)
  =
  \eta\log\frac{4\eta}{(1+\eta)^2}
  +
  R\log\frac{2}{1+\eta}.
\end{equation}
Subtracting gives
\begin{align}
  \left|
  \Pi_{U_1}^{\sx}(p)-\Pi_{U_1}^{\sx}(q)
  \right|
  &=
  2\eta\log\frac{1+\eta}{2\eta}
  \notag\\
  &\quad+
  (1-\eta)\log\frac{1+\eta}{1-\eta}.
\end{align}
Both terms are positive for \(0<\eta<1\).  Dividing by
\(\eta\log(1/\eta)\) and taking \(\eta\downarrow0\) gives
\begin{equation}
  \lim_{\eta\downarrow0}
  \frac{
    \left|\Pi_{U_1}^{\sx}(p)-\Pi_{U_1}^{\sx}(q)\right|
  }{
    \eta\log(1/\eta)
  }
  =2.
\end{equation}
This gives one fixed system and atom with net coefficient two.
\end{proof}

\begin{remark}[Sharpness scope]
Proposition~\ref{prop:leading-sharp} establishes only the common worst-case leading coefficients
for families of fixed-system bounds whose \(O_{\mathcal F}(\eta)\) remainder may depend on the
system.  It does not claim that every system or atom attains either coefficient, and it does not
make the lower-order terms, branch terms, or every individual cumulative and atom bound globally
sharp.
\end{remark}

\section{Formal and executable evidence boundary}

The Lean artifact checks finite-vector overlap and residual algebra, residual-entropy ceilings,
abstract component and signed transfer lemmas, and a generic down-set M{\"o}bius row-sum identity.
It also defines the exact heterogeneous dependent Cartesian key, the keyed
shared-exclusions events \(A_\beta\), \(B_\beta\), and \(C\), their equivalence-class covers,
anchors, intersection identity, positive event masses, and law-independent event map.  A separate
checked module proves the support-restricted finite load bound \(0\leq T_N\leq J\eta\), including
its concrete \(A_\beta\), \(B_\beta\), and \(C\) corollaries.  Lean also checks endpoint and linear
properties of \(g_J\).  A separate checked bridge starts from supplied exact natural counts with
positive total and proves event-count, positivity, local-logarithm, and positive-support average
identities for the source-one, source-two, joint-source, and redundancy signed-net cumulative
nodes.

Lean does not yet bundle two-law semantic total variation, derive the support-change \(g_J\)
transfer from \(T_N\), identify concrete full-lattice M{\"o}bius atoms, import the published
pointwise sign theorem, or prove the complete informative or misinformative averaged component
family.  Bytes or rows to counts, sampling and population claims, Rust or floating-point
refinement, and certifier/parser execution also remain outside the checked theorem.

The implementation-separated generator is standard-library-only and imports no pid-rs code.  It
retains rational laws, exact integer identities, high-precision logarithmic references, equality
witnesses, and a bounded seeded challenge corpus.  The larger historical grid searches remain
separate revision-1 route evidence.  This separation reduces shared implementation risk; it is not
independent authorship or review.  Decimal comparison is not certified interval arithmetic.  The
Rust comparison replays committed empirical count tables.  It is not a universal refinement proof
or a portable binary64 enclosure.

The dependency-colored probability theorem has a separate premise: complete rows in each declared
color must be mutually independent.  The code cannot infer that premise from a finite sample.

\section{Release interpretation}

The exact-real mathematical result supports this statement:
\begin{quote}
Averaged categorical shared-exclusions cumulatives and full-lattice atoms are quantitatively
continuous on a fixed finite closed simplex, including across support changes, without a positive
minimum cell-mass premise.
\end{quote}

It does not support these statements:
\begin{itemize}[leftmargin=2em]
  \item every pointwise atom has a finite support-boundary limit;
  \item the functional is globally Lipschitz;
  \item the result is independent of alphabet size;
  \item a binary64 atom sign is certified;
  \item a dependence coloring is valid because autocorrelation is small;
  \item continuous PID2 is generally consistent;
  \item full continuous PID3 is ready for scientific inference; or
  \item a PID value can independently authorize an action.
\end{itemize}

\begin{thebibliography}{9}

\bibitem{makkeh2021}
A.~Makkeh, A.~J. Gutknecht, and M.~Wibral.
\newblock Introducing a differentiable measure of pointwise shared information.
\newblock \emph{Physical Review E}, 103:032149, 2021.
\newblock \url{https://doi.org/10.1103/PhysRevE.103.032149}.

\bibitem{gutknecht2021}
A.~J. Gutknecht, M.~Wibral, and A.~Makkeh.
\newblock Bits and pieces: understanding information decomposition from part-whole relationships
and formal logic.
\newblock \emph{Proceedings of the Royal Society A}, 477:20210110, 2021.
\newblock \url{https://doi.org/10.1098/rspa.2021.0110}.

\bibitem{audenaert2007}
K.~M.~R. Audenaert.
\newblock A sharp continuity estimate for the von Neumann entropy.
\newblock \emph{Journal of Physics A}, 40:8127--8136, 2007.
\newblock \url{https://doi.org/10.1088/1751-8113/40/28/S18}.

\end{thebibliography}

\end{document}
